\chapter[TUNNELING RÉSEAU ET DIGIPEATING TIERS]{Tunneling réseau et digipeating
tiers}\label{tunneling-ruxe9seau-et-digipeating-tiers}

\section{Réseaux tiers}\label{ruxe9seaux-tiers}

APRS fournit un mécanisme pour tunneliser des paquets à travers des
réseaux non-AX.25 (Internet, LAN Ethernet, câble direct). Ces réseaux ne
comprennent pas les adresses AX.25, il faut les transmettre comme
données avec la charge utile originale.

\section{Source Path Header}\label{source-path-header}

Avant d'envoyer un paquet dans le réseau tiers, on préfixe l'adresse
APRS au DTI et au reste de la data. C'est le \textbf{Source Path Header}
--- source, destination et digipeaters avec SSIDs. But : permettre au
récepteur distant d'identifier l'émetteur (nécessaire pour ACK, utile
pour diagnostic).

Deux formats en usage : \textbf{TNC-2} et \textbf{AEA}. L'APRS Working
Group a acté la standardisation sur TNC-2 pour l'avenir. Les TNCs AEA
produisent du TNC-2 quand \texttt{BBSMSGS} est ON.

\textbf{TNC-2} --- astérisque après le digi qui a effectivement entendu
:

\begin{verbatim}
Source(-SSID) > Destination(-SSID) , [Digi(-SSID)(*),]... :
\end{verbatim}

Ex : \texttt{WB4APR-14\textgreater{}APRS,RELAY*,WIDE:} (le WIDE est «
inutilisé »).

\textbf{AEA} --- astérisque après la source ou le digi qui a
effectivement entendu :

\begin{verbatim}
Source(-SSID)(*) > [Digi(-SSID)(*)>]... Destination(-SSID) :
\end{verbatim}

Ex :
\texttt{WB4APR-14\textgreater{}RELAY*\textgreater{}WIDE\textgreater{}APRS:}.

Dans les deux formats, SSID omis si \texttt{-0}. Les digipeaters après
l'astérisque sont dits « unused » --- supprimés lors de la construction
du Third-Party Header.

\section{Third-Party Header}\label{third-party-header}

Après émergence d'un paquet du réseau tiers, la gateway réceptrice
insère un DTI \texttt{\}} et modifie le Source Path Header pour former
le \textbf{Third-Party Header}, avant retransmission sur l'APRS local.

\begin{verbatim}
} | Third-Party Header | Rest of original data
\end{verbatim}

\textbf{TNC-2} :

\begin{verbatim}
Source Path Header (sans digi unused, sans *, sans :) , Third-Party Network ID , Receiving Gateway(-SSID) * :
\end{verbatim}

Ex : \texttt{\}WB4APR-14\textgreater{}APRS,RELAY,TCPIP,G9RXG*:...}

Suppression des digis inutilisés et de l'astérisque. Puis insertion de 2
indicatifs supplémentaires : Third-Party Network Identifier (ex :
\texttt{TCPIP}) et indicatif de la gateway réceptrice suivi d'un
\texttt{*}.

\section{Action à la réception}\label{action-uxe0-la-ruxe9ception}

La station peut extraire l'indicatif de l'émetteur d'origine (pour ACK)
et connaît les autres adresses (diagnostic réseau).

\section{Exemple : message via
Internet}\label{exemple-message-via-internet}

Scénario : - WB4APR-14 veut envoyer à G3NRW via Internet. - Gateway la
plus proche de WB4APR-14 : K4HG (accessible via \texttt{RELAY,WIDE}). -
Gateway la plus proche de G3NRW : G9RXG.

Processus : 1. WB4APR-14 construit \texttt{:G3NRW\ \ \ \ :Hi\ Ian\{001}
et le transmet via UNPROTO \texttt{RELAY,WIDE}. 2. K4HG reçoit via RELAY
et construit
\texttt{WB4APR-14\textgreater{}APRS,RELAY*,WIDE::G3NRW\ \ \ \ :Hi\ Ian\{001}.
3. K4HG envoie ce paquet en telnet à un APRServer. 4. Toutes les
gateways Internet APRServe (dont G9RXG) reçoivent le paquet. 5. G9RXG
convertit :
\texttt{\}WB4APR-14\textgreater{}APRS,RELAY,TCPIP,G9RXG*::G3NRW\ \ \ \ :Hi\ Ian\{001}
(le WIDE inutilisé est supprimé). 6. G9RXG transmet sur l'APRS local. 7.
G3NRW reçoit, strippe le Third-Party Header, décode le message pour lui,
ACK routé vers WB4APR-14.

